% GENERATED FILE. Edit canonical lessons, then run npm run generate:books.
% canonical-source-hash: fnv1a64:468293b96eb9da1a
% canonical-lessons: PT-C385-burocracia, PT-C385-processo, PT-C385-pedido, PT-C385-legal, PT-C385-ilegal

\chapter{A burocracia}
\label{ch:pt-a-burocracia}
\hlchaptermodality{385}

\begin{chapteropening}
\textbf{By the end of this chapter:} \emph{I can deal with red tape: the process, the request, and what is legal or illegal.}

Complete the last lesson of chapter 385: I can deal with red tape: the process, the request, and what is legal or illegal.
\end{chapteropening}

\section[a burocracia]{a burocracia — the bureaucracy, the red tape}

\label{lesson:PT-C385-burocracia}



\begin{quote}
Before the new one: say the Portuguese for to renew, then the Portuguese for the certificate.
\end{quote}

\subsection*{You'll want to know: a burocracia}
\textbf{a burocracia} — \textquotedblleft{}the bureaucracy, the red tape\textquotedblright{}.

\begin{quote}
Para abrir uma empresa há muita burocracia.
\end{quote}

\begin{quote}
\emph{There is a lot of red tape to open a company.}
\end{quote}

\emph{Burocrático} — bureaucratic.

\begin{cousinweb}
\emph{burocracia} is usually traced to French \emph{bureaucratie}, which EtymDB builds from \emph{bureau} \textquotedblleft{}office, desk\textquotedblright{} and \emph{-cratie} \textquotedblleft{}rule of\textquotedblright{} (EtymDB); \emph{-cratie} goes back to Greek \emph{krátos} \textquotedblleft{}power\textquotedblright{}.
\end{cousinweb}

\begin{grammarlens}[title={What you've built}]
One more word for telling, explaining and describing at B1.
\end{grammarlens}

\subsection*{Your turn}
\begin{itemize}
\raggedright
  \item \emph{Say it:} \emph{a burocracia}
  \item \emph{Say it:} \emph{a burocracia}, once more
  \item \emph{Say it:} a sentence of your own with \emph{a burocracia}
  \item \emph{Recall:} say the Portuguese for to renew, then the Portuguese for the certificate, then say \emph{a burocracia} again
\end{itemize}

\begin{tcolorbox}[breakable,colback=teal!4,colframe=teal!35!black,title={Before you move on}]
What does a burocracia mean? (\textquotedblleft{}The bureaucracy, the red tape\textquotedblright{}.) Say it once more.
\end{tcolorbox}
\section[o processo]{o processo — the process; the case, the lawsuit}

\label{lesson:PT-C385-processo}



\begin{quote}
Before the new one: say the Portuguese for the certificate, then the Portuguese for the bureaucracy.
\end{quote}

\subsection*{You'll want to know: o processo}
\textbf{o processo} — \textquotedblleft{}the process; the case, the lawsuit\textquotedblright{}.

\begin{quote}
O processo demorou três meses.
\end{quote}

\begin{quote}
\emph{The process took three months.}
\end{quote}

In law, \emph{um processo} is a case: \emph{pôr um processo a alguém} — to sue someone.

\begin{cousinweb}
\emph{processo} is a learned borrowing from Latin \emph{prōcessus} \textquotedblleft{}a going forward, an advance\textquotedblright{} (EtymDB), from \emph{prōcēdere}, built on \emph{cēdere} \textquotedblleft{}to go\textquotedblright{} (Lewis \& Short).
\end{cousinweb}

\begin{grammarlens}[title={What you've built}]
One more word for telling, explaining and describing at B1.
\end{grammarlens}

\subsection*{Your turn}
\begin{itemize}
\raggedright
  \item \emph{Say it:} \emph{o processo}
  \item \emph{Say it:} \emph{o processo}, once more
  \item \emph{Say it:} a sentence of your own with \emph{o processo}
  \item \emph{Recall:} say the Portuguese for the certificate, then the Portuguese for the bureaucracy, then say \emph{o processo} again
\end{itemize}

\begin{tcolorbox}[breakable,colback=teal!4,colframe=teal!35!black,title={Before you move on}]
What does o processo mean? (\textquotedblleft{}The process; the case, the lawsuit\textquotedblright{}.) Say it once more.
\end{tcolorbox}
\section[o pedido]{o pedido — the request, the application; the order (in a restaurant)}

\label{lesson:PT-C385-pedido}



\begin{quote}
Before the new one: say the Portuguese for the bureaucracy, then the Portuguese for the process.
\end{quote}

\subsection*{You'll want to know: o pedido}
\textbf{o pedido} — \textquotedblleft{}the request, the application; the order (in a restaurant)\textquotedblright{}.

\begin{quote}
Fiz um pedido à câmara para abrir uma esplanada.
\end{quote}

\begin{quote}
\emph{I made an application to the council to open a terrace café.}
\end{quote}

\emph{Fazer um pedido} — to make a request; \emph{a pedido de} — at the request of.

\begin{cousinweb}
\emph{pedido} is the participle of \emph{pedir} used as a noun, from Latin \emph{petītus} \textquotedblleft{}sought\textquotedblright{} (EtymDB); \emph{pedir} is inherited from Latin \emph{petere} \textquotedblleft{}to seek, to ask\textquotedblright{} (Lewis \& Short).
\end{cousinweb}

\begin{grammarlens}[title={What you've built}]
One more word for telling, explaining and describing at B1.
\end{grammarlens}

\subsection*{Your turn}
\begin{itemize}
\raggedright
  \item \emph{Say it:} \emph{o pedido}
  \item \emph{Say it:} \emph{o pedido}, once more
  \item \emph{Say it:} a sentence of your own with \emph{o pedido}
  \item \emph{Recall:} say the Portuguese for the bureaucracy, then the Portuguese for the process, then say \emph{o pedido} again
\end{itemize}

\begin{tcolorbox}[breakable,colback=teal!4,colframe=teal!35!black,title={Before you move on}]
What does o pedido mean? (\textquotedblleft{}The request, the application; the order (in a restaurant)\textquotedblright{}.) Say it once more.
\end{tcolorbox}
\section[legal]{legal — legal, lawful; cool (informal, in Brazil)}

\label{lesson:PT-C385-legal}



\begin{quote}
Before the new one: say the Portuguese for the process, then the Portuguese for the request.
\end{quote}

\subsection*{You'll want to know: legal}
\textbf{legal} — \textquotedblleft{}legal, lawful; cool (informal, in Brazil)\textquotedblright{}.

\begin{quote}
Fazer isso é legal, mas não é bonito.
\end{quote}

\begin{quote}
\emph{Doing that is legal, but it is not nice.}
\end{quote}

In Brazil, informal \emph{legal!} means cool, great; in Portugal it is mostly lawful.

\begin{cousinweb}
\emph{legal} is a learned borrowing from Latin \emph{lēgālis} \textquotedblleft{}legal\textquotedblright{} (EtymDB), built on \emph{lēx, lēgis} \textquotedblleft{}law\textquotedblright{} (Lewis \& Short); the same Latin word gave the inherited \emph{leal} \textquotedblleft{}loyal\textquotedblright{}.
\end{cousinweb}

\begin{grammarlens}[title={What you've built}]
One more word for telling, explaining and describing at B1.
\end{grammarlens}

\subsection*{Your turn}
\begin{itemize}
\raggedright
  \item \emph{Say it:} \emph{legal}
  \item \emph{Say it:} \emph{legal}, once more
  \item \emph{Say it:} a sentence of your own with \emph{legal}
  \item \emph{Recall:} say the Portuguese for the process, then the Portuguese for the request, then say \emph{legal} again
\end{itemize}

\begin{tcolorbox}[breakable,colback=teal!4,colframe=teal!35!black,title={Before you move on}]
What does legal mean? (\textquotedblleft{}Legal, lawful; cool (informal, in Brazil)\textquotedblright{}.) Say it once more.
\end{tcolorbox}
\section[ilegal]{ilegal — illegal}

\label{lesson:PT-C385-ilegal}



\begin{quote}
Before the new one: say the Portuguese for the request, then the Portuguese for legal.
\end{quote}

\subsection*{You'll want to know: ilegal}
\textbf{ilegal} — \textquotedblleft{}illegal\textquotedblright{}.

\begin{quote}
Estacionar em cima do passeio é ilegal.
\end{quote}

\begin{quote}
\emph{Parking on the pavement is illegal.}
\end{quote}

The opposite of \emph{legal}.

\begin{cousinweb}
\emph{ilegal} is a learned borrowing from Medieval Latin \emph{illēgālis} (EtymDB), built from \emph{in-} \textquotedblleft{}not\textquotedblright{} and \emph{lēgālis} \textquotedblleft{}legal\textquotedblright{}, on \emph{lēx} \textquotedblleft{}law\textquotedblright{} (Lewis \& Short).
\end{cousinweb}

\begin{grammarlens}[title={What you've built}]
One more word for telling, explaining and describing at B1.
\end{grammarlens}

\subsection*{Your turn}
\begin{itemize}
\raggedright
  \item \emph{Say it:} \emph{ilegal}
  \item \emph{Say it:} \emph{ilegal}, once more
  \item \emph{Say it:} a sentence of your own with \emph{ilegal}
  \item \emph{Recall:} say the Portuguese for the request, then the Portuguese for legal, then say \emph{ilegal} again
\end{itemize}

\begin{tcolorbox}[breakable,colback=teal!4,colframe=teal!35!black,title={Before you move on}]
What does ilegal mean? (\textquotedblleft{}Illegal\textquotedblright{}.) Say it once more.
\end{tcolorbox}
